\histitem{Fredholm index and perturbations}

The notion of index appears in 1920 in a work by Fritz Noether on integral equations. In his lecture at the 1904 International Congress, Hilbert considered a problem posed by Riemann in potential theory: given a bounded open domain \(\Omega\) of the plane whose boundary is a smooth simple closed curve \(\Gamma\) , as well as three functions \(a\) , \(b\) , \(f\) on \(\Gamma\) , the problem is to find a function \(z: \overline{\Omega} \to \mathbf{C}\) , holomorphic on \(\Omega\) and continuous on \(\overline{\Omega}\) , satisfying

\[
a \mathcal {R} (z) + b \mathcal {I} (z) = f \quad \text { on } \Gamma .
\]

Parametrizing \(\Gamma\) by a real number \(s \in [0, 2\pi]\) , Hilbert showed that the problem reduces to solving an integral equation “with singular kernel” for \(\varphi = \mathcal{R}(z)\) , namely

\[
a (s) \varphi (s) + \int_ {0} ^ {2 \pi} \mathrm{K} (s, t) \varphi (t) d t = f (s).\tag{4}
\]

Here the kernel \(\mathrm{K}(s, t)\) takes the form \(b(s) \cot\left(\frac{t-s}{2}\right) + \mathrm{A}(s, t)\) for a continuous function A. It is therefore singular along the diagonal, and the above integral is to be taken in the sense of the Cauchy principal value.

Noether {[}53{]} observes that, unlike integral equations satisfying the Fredholm alternative, in the case of a kernel K as above and a nonzero right-hand side \(f\), it is possible that (4) admits families of nontrivial solutions. He shows that the space of solutions is governed by the integer

\[
n = \frac {1}{2 \pi} \int_ {\Gamma} d (\log (a - i b)).\tag{5}
\]

If \(n < 0\), there is no nontrivial solution, and if \(n \geqslant 0\), equation (4) admits a family of solutions depending on 2n parameters. Noether uses the term "index" (Index) for the integer n and recognizes in it a winding number, a notion classical since the end of the nineteenth century in the study of functions of a complex variable.

Later works would go on to study examples of equations requiring an extension of Fredholm theory. But it would take more than twenty years before the notion of a Fredholm map was systematically extracted, at the same time as the study of perturbations by compact operators matured.

In 1909, H. Weyl {[}81{]} had shown that if two quadratic forms bounded on \(\ell_{\mathbf{C}}^{2}(\mathbf{N})\) have a completely continuous difference, then their essential spectra coincide (cf. III, p. 89, th. 2). Moreover, Riesz had noted in his 1916 memoir, naturally without the modern language, that the compact operators on a Banach space E form a two-sided ideal of \(\mathcal{L}(\mathrm{E})\). It seems, however, that the path indicated by these two remarks was not explored until the 1940s. In 1941, J. Calkin, motivated by the work of J. von Neumann on operator algebras (which we shall have occasion to discuss shortly, see p. 537), pointed out the interest of studying congruences modulo this ideal in the framework of Hilbert spaces {[}10{]}. He shows how the structure of the algebra that now bears his name (cf. III, p. 75) makes it possible to recover Weyl's result simply.

Meanwhile, around 1942, J. Dieudonné {[}15{]} became interested in the idea of studying perturbations of a continuous linear map \(u\) between normed spaces, in the case where the perturbation is “small” in the sense of the operator norm. He restricts himself to the case where \(u\) is a strict morphism whose kernel is finite-dimensional or of finite codimension; he proves that the same then holds for any “small perturbation” of \(u\) , and shows that in the case of Fredholm maps ( \(cf. n^{\circ} 2\) from III, p. 40), the index is locally constant (th. 1 of III, p. 58), as are several of the results presented in § 4 of III, p. 55.

These two ideas converge around 1950, when F. Atkinson {[}3{]}, I. Gohberg {[}26{]}, S. Mikhlin {[}46{]} , and B. Yood {[}93{]} study perturbations by compact operators. They explicitly make the connection with the Noether index and bring out the notions of Fredholm map {[}3, p. 8{]} and of Riesz map {[}93, §5{]}. The latter are the subject of systematic research in the following decade, where the results presented in Chapter III take on more or less their current form.

The usual framework for these theories is that of Banach spaces; however, various applications motivate their generalization to broader classes of topological vector spaces. Thus, the case of Fréchet spaces appears naturally in 1954 in works of H. Cartan and J-P. Serre {[}13{]} which demonstrate the finiteness of the cohomology of a compact complex analytic manifold with coefficients in a coherent sheaf, the proof invoking deformation results proved by L. Schwartz {[}66{]} (cf. th. 2 of III, p. 73).

The invariance of the index under deformation is perfectly illuminated, in the case of formula (5), if one observes that the right-hand side (winding number) is invariant under homotopy. Such topological expressions of the index would soon become famous for Fredholm applications arising from differential geometry. If D is a differential operator on a compact \(C^{\infty}\) manifold X (cf. VAR, §14), and if D is “elliptic” (cf.{[}2, §1{]}), then an extension of D to suitable Sobolev spaces defines a Fredholm map. The search for a formula giving a topological expression for its index using the “characteristic classes” of X, suggested among others by the work of I. Gelfand around 1960 {[}23, vol. I, p. 65{]}, will lead in 1963 to the “index theorem” of M. Atiyah and I. Singer {[}2{]}. This latter would give rise to extraordinary developments at the confluence of analysis, topology, and differential geometry. We cannot describe here the vast fields they opened, which, still very fertile today, have not ceased to draw nourishment from Fredholm theory.

\histitem{Partial operators and the spectral theorem}

General Hilbertian spectral theory was created in large part to address the problem of the mathematical foundations of quantum mechanics, notably under the impetus of von Neumann.

We refer to the work of M. Jammer {[}35{]} for a detailed description of the development of quantum theory up to the fundamental paper by W. Heisenberg that introduced quantum mechanics {[}28{]}. Less than five years elapsed between its publication in the summer of 1924 and that of the article {[}50{]} in 1929, where von Neumann presents in a perfectly lucid and rigorous manner all the fundamental results of the theory of partial (often called "unbounded") Hermitian operators on Hilbert spaces.

By giving the first axiomatic presentation of Hilbert spaces {[}49{]}, von Neumann crowned an idea long in gestation (cf. ÉHM, p. 267), and illuminated various isomorphisms between spaces of sequences or functions well known to Schmidt, Fréchet, Fischer, and Riesz before 1910. This rise in abstraction allowed von Neumann to make considerable conceptual advances.

Such is the case with the idea of a partial operator, which seems to hover, still formless, over several of the works that, after Hilbert, took up Fredholm equations again (see above, no. 1). Hilbert had noticed that many differential equations of Sturm--Liouville type, in particular in the case of an unbounded interval, could be reduced to Fredholm equations; but to apply the integral methods, it was necessary to weaken the conditions imposed on the kernel and to leave the framework in which the results of Fredholm, Hilbert, and Schmidt could be applied. Over the years, more and more singular kernels were studied, until equations were considered for which the left-hand side makes sense only if the unknown belongs to a subspace of L²(I). In this direction, one must mention works of E. Hilb {[}30{]} in 1908, of H. Weyl on Sturm--Liouville theory {[}82{]} in 1910, and above all of T. Carleman {[}11{]} in 1923, in a very general framework.

By unifying these problems, von Neumann cast entirely new light on these classical works. He emphasized in particular the fundamental distinction between symmetric and self-adjoint operators, and interpreted the self-adjoint extensions of a symmetric operator in terms of abstract "boundary conditions," generalizing the well-known boundary conditions that appeared, for example, in Dirichlet and Neumann problems for the Laplace operator on a bounded open set of \(R^{n}\) .

At the same time, von Neumann {[}51{]} brought to the fore the notion of normal operator \(^{(5)}\) as the natural framework for spectral theory, insisting on the role of the commutative algebra generated by the operator and its adjoint.

Relying on this precise mathematical formalism, which also revealed the equivalence between the viewpoints of Heisenberg and Schrödinger, and on the Copenhagen interpretation of experimental procedures and results, non-relativistic quantum mechanics from then on reached a state of refinement that has hardly evolved since; the validity of this physical theory, despite its surprising, even paradoxical-looking consequences, has since been constantly confirmed by experiment, and with remarkable precision.

From a formal point of view, the foundations of spectral theory were laid by von Neumann. In parallel, after the presentation of the latter's first ideas, M. Stone had pursued similar research between 1928 and 1930 and obtained closely related results. The clear and complete presentation of the known results that Stone published in 1932 {[}72{]} had an important influence on the diffusion of Hilbertian spectral theory.

Nevertheless, other improvements in the presentation have been important for the diffusion of these results. Indeed, the spectral theorem (cf. th. 1 of IV, p. 266), as presented by von Neumann, was long considered very difficult, largely because it was stated in the framework of vector-valued measures.

Later, P. Halmos {[}27{]}, by emphasizing the interpretation of the spectral theorem by means of multiplication operators on \(L^{2}\) spaces, revealed its essentially elementary aspect in the case of bounded operators. Quite recently, S. Woronowicz (in the somewhat different framework of regular operators on Hilbert modules over C*-algebras {[}92{]}) brought a simplification by introducing

the "bornification" (cf. no. 2 of IV, p. 265), which allows one to treat normal partial operators as simply as self-adjoint operators.

Moreover, if D is a formally symmetric scalar differential operator on an open set U of \(R^{n}\) , the study of the self-adjoint extensions of D to \(L^{2}(U)\) is naturally linked to the study of distributions on U (in the case of the Laplacian, cf. no. 6 of IV, p. 242). We refer to the work of J. Dieudonné {[}16, ch. VII{]} for the history of the latter --- a sinuous history that unravels shortly after the period we have just evoked, with decisive contributions of S. Sobolev {[}69{]} in 1936, and of L. Schwartz {[}65{]} ten years later.

From a purely mathematical point of view, spectral theory interacts in a particularly fruitful way with Riemannian geometry. Thus, taking up a question of S. Bochner, M. Kac gave in 1966 a strikingly vivid presentation {[}36{]} of the problem of determining which geometric properties of a Riemannian manifold can be deduced from the spectral properties of the Laplace operator. This question had been anticipated by the physicist A. Schuster {[}64{]} in 1882: To find out the different tunes sent out by a vibrating system is a problem which may or may not be solvable in certain special cases, but it would baffle the most skillful mathematician to solve the inverse problem and to find out the shape of a bell by means of the sounds which it is capable of giving out. (“Determining the harmonics emitted by a vibrating system is a problem which may or may not be solvable in certain special cases, but the most skillful mathematician would be baffled if he had to solve the inverse problem and determine the shape of a bell from the sounds it can emit”).

\histitem{Connection between harmonic analysis and group theory}

We have described, in the historical notes of the Book on Integration (cf. ÉHM, p. 275), how questions related to the heat equation led Fourier to the revolutionary conception of representing an arbitrary function as a sum of trigonometric functions.

We cannot trace here the development of the theory of Fourier series and integrals throughout the nineteenth century, nor the profound influence that the questions raised by this theory had on the evolution of Analysis --- to the point of overturning the conception of real numbers, contributing to the birth of set theory (cf. ÉHM, p. 42). But in order to understand the form of the ideas set forth in this Book, it is appropriate to describe the manner in which harmonic analysis was connected after 1925 to group theory and to Hilbertian spectral theory.

Recognition of the link between Fourier's ideas and the group structures of R/Z and R, as well as the generalization of these ideas to other groups, came rather late.

One can naturally, in the early stages of the theory of finite groups and their characters, identify results that today appear to contain the rudiments of finite Fourier analysis. For example, the dual of the group of invertible elements of Z/qZ, as well as the orthogonality relation expressing the characteristic function of an element as a linear combination of characters, appear implicitly in 1837 in Dirichlet's article {[}17{]} proving that there are infinitely many prime numbers \(p \equiv a \mod q\) if a is relatively prime to q. He relies on ideas of Gauss, who had introduced the term “character” in his study of binary quadratic forms with integer coefficients and fixed discriminant (cf. {[}22, § 230{]}).

But neither Gauss nor Dirichlet uses the language of groups. It was Dedekind, presenting these works in 1879, who recognized their latent role and defined the notion of character for commutative groups. Dedekind's remarks were considerably amplified by H. Weber, first in 1882 {[}77{]}, then in 1886 {[}78, p. 112{]}, where he introduced the dual group of a finite abelian group A and noted that the latter is isomorphic to A. His purpose was then to construct abelian extensions of the field Q with given Galois group, and he did not consider the problem of harmonic analysis on such a group.

In the case of noncommutative groups, the foundations of the theory of representations of finite groups were solidly established around the turn of the twentieth century, following the works of G. Frobenius, W. Burnside, and I. Schur (cf. ÉHM, p. 154). After 1905, it was clear that the orthogonality relations of characters played a crucial role in the organization of the theory. But their kinship with harmonic analysis remained hidden.

It is to Hermann Weyl that we owe the connection between these fields, when he conceived the idea of a general theory of representations for compact groups. In several major texts published between 1925 and 1927, he recognized the links that would unite such a theory with Fourier analysis and Hilbertian spectral theory, and he laid the foundations for many later discoveries.

It was a letter from Schur that set Weyl on this path, thanks to their common interest in invariant theory. We described, in the note on Haar measure (ÉHM, pp. 289–291), how A. Hurwitz had conceived as early as 1897 the idea of using an “invariant integration” to determine the polynomials on \(R^{n}\) invariants under the orthogonal group. It seems that Schur was not aware of Hurwitz's results before 1924, when he suddenly understood how invariant integration could be used to extend to the group \(\mathbf{O}(n)\) character theory and the orthogonality relations he had established for finite groups. Weyl immediately saw that Schur's methods, combined with the work of Élie Cartan, made it possible to construct the irreducible representations of compact semisimple Lie groups. The series of three memoirs in which Weyl blended the ideas of Cartan and Schur, published in 1925 {[}83{]}, already contained the essential results of LIE, IX, §6–7; cf. ÉHM, p. 328–330.

But it was the following year, in an equally famous text written with his student F. Peter {[}54{]}, that Weyl brought to light the connections between the representation theory of compact groups, Fourier analysis, and Hilbertian spectral theory. This text of fewer than twenty pages contains in embryo a great many future developments.

Peter and Weyl free themselves of any algebraic hypothesis on the structure of the group under study, assuming only that G is a compact topological group equipped with an invariant measure. The existence of such a measure was clear for connected Lie groups; soon A. Haar would establish it in general, partly motivated by the results presented here (cf. ÉHM, p. 291).

The starting point of their study is the orthogonality of matrix coefficients. Schur had observed that by choosing a representative \(\pi \in \widehat{\mathrm{G}}\) for each equivalence class of irreducible representations, a G-invariant inner product on the space E of \(\pi\) , and an orthonormal basis of E, one obtains a family of matrix coefficients that is orthonormal in \(\mathrm{L}^2 (\mathrm{G})\) . Peter and Weyl intend to show that such a family is complete.

In the case of finite groups, the analogous result had been proved by Frobenius by studying the group algebra C{[}G{]}. Peter and Weyl then consider the space \(\mathcal{C}(\mathrm{G})\) of continuous functions on G and equip it with the convolution product. This appears to be the first use of the convolution product as an abstract operation in explicit connection with the group structure (cf. ÉHM, p. 295). Peter and Weyl moreover call Gruppenzahl ("group number") an element of the algebra \(\mathcal{C}(\mathrm{G})\) , and denote by xy the (convolution) product of two Gruppenzahlen. They also equip \(\mathcal{C}(\mathrm{G})\) with the involution \(f \mapsto \widetilde{f}\) , where \(\widetilde{f}(g) = \overline{f(g^{-1})}\) for \(g \in G\) .

Having thus introduced the involutive algebra \(\mathcal{C}(\mathrm{G})\) , the first observation made by Peter and Weyl is that every unitary representation \(\pi\) of G gives rise to a representation \(f \mapsto \pi(f)\) of \(\mathcal{C}(\mathrm{G})\) (cf. V, p. 400). The notion of Hermitian element of \(\mathcal{C}(\mathrm{G})\) is explicitly defined, as is (implicitly) that of positive element, which opens the way to the application of Hilbert space techniques.

Peter and Weyl no longer hesitate to call the operator \(\pi(f)\) the "Fourier coefficient" of \(f\) , and continue the parallel with Fourier theory by noting that Schur's orthogonality relations imply Bessel's inequality

\[
\sum_ {\pi \in \widehat {\mathrm{G}}} \dim (\pi) \operatorname{Tr} (\pi (f) \pi (f) ^ {*}) \leqslant (f * \widetilde {f}) (e) = \| f \| _ {2} ^ {2}.\tag{6}
\]

The Hilbert--Schmidt spectral theory then allows them to show that this is an equality ("Plancherel formula"). To every element \(\varphi\) of \(\mathcal{C}(\mathrm{G})\) , Peter and Weyl associate the continuous kernel \(K_{\varphi}\colon G\times G\to C\) defined by \(\mathrm{K}_{\varphi}(x,y)=\varphi(xy^{-1})\) . They then observe that if \(\varphi\) is of the form \(f*\widetilde{f}\) with \(f\in\mathcal{C}(\mathrm{G})\) , \(^{(6)}\) then \(K_{\varphi}\) is a positive Hermitian kernel in the sense of Hilbert--Schmidt theory. A variant of a Schmidt algorithm \(^{(7)}\) then allows them to construct an orthonormal basis of eigenfunctions of the kernel \(K_{\varphi}\) , and then to reduce the proof of the equality in (6) to the fact that the trace of the operator defined by \(K_{\varphi}\) is equal to the sum of its eigenvalues.

Naturally, only the representations \(\pi\) satisfying \(\pi(f) \neq 0\) can occur in (6); the spectral theory for \(K_{\varphi}\) shows the existence of such a representation if \(f\) is not identically zero. Peter and Weyl easily deduce from this that the irreducible representations separate the points of G --- a special case of the Gelfand--Raikov theorem, which will be discussed later.

In the celebrated works of Frobenius, an equality analogous to (6) made it possible to show that the regular representation of G contains all irreducible representations. But this involved applying this equality by taking for f the identity element of C{[}G{]}; one then obviously had \(\pi(f) \neq 0\) for all \(\pi\) . According to Peter and Weyl, the absence of an identity element for convolution explains many difficulties in their proof; they nevertheless observe, borrowing from the theory of Fourier series an idea destined for a great future, that one can deduce the analogue of Frobenius's result by applying (6) to a sequence of functions forming an approximate identity for convolution (cf. no. 10 of I, p. 120).

Likewise, starting from the equality arising from (6), the polarized version

\[
\sum_ {\pi \in \widehat {\mathrm{G}}} \dim (\pi) \operatorname{Tr} (\pi (x) \pi (y)) = (x * y) (e)\tag{7}
\]

makes it possible to obtain, by taking for \(y\) an approximate identity for convolution, a uniform approximation of every element of \(\mathcal{C}(\mathrm{G})\) by linear combinations of matrix coefficients of irreducible representations. These methods had just enabled Weyl to give a new proof {[}84{]} of the fundamental result of the theory of "almost periodic functions" of H. Bohr, to which we shall return shortly. Peter and Weyl point out that one can read therein the first analytic application of the theory of representations of a noncompact group, here that of the translations of \(\mathbf{R}\) .

The young quantum theory very quickly provided Weyl with the opportunity to return to the representations of R. At the end of 1927, he noted the interest of groups (finite or not) and of representations for clarifying the foundations of this theory {[}85{]}. He would return to this the following year, in a resounding book {[}86{]}.

In his 1927 article, Weyl observes that if u is a continuous self-adjoint operator on a Hilbert space, then \(t \mapsto e^{itu}\) is a unitary representation of R. But there exist unitary representations of R that are not of this form; taking up von Neumann's idea on the role of symmetric partial operators in representing observable physical quantities, Weyl suggests that every unitary representation of R is of the form \(t \mapsto e^{itu}\) where u is a partial self-adjoint operator on a Hilbert space. This point of view allows him to reduce the analysis of the "canonical relations" of Heisenberg and Schrödinger, already studied in detail by von Neumann, to that of the connections between two unitary representations of R on the same Hilbert space. Stone {[}73{]} shows in 1932 the result hoped for by Weyl (V, p. 428, th. 1), in the wake of his study of the spectral properties of self-adjoint operators.

\histitem{Locally compact commutative groups}

The interaction between harmonic analysis and group representations was thus firmly established by the early 1930s. It was greatly strengthened in the following decade, with each of the two subjects prompting rapid progress in the other.

A very fertile ground for this interaction was the study of almost periodic functions, introduced by H. Bohr in 1924 in the case of functions of a real variable {[}9{]}. By one of the fundamental theorems of the latter, these are the bounded continuous functions that are uniform limits on \(\mathbf{R}\) of linear combinations of functions \(x \mapsto e^{i\lambda x}\) , \(\lambda \in \mathbf{R}\) (cf. II, p. 292, exerc. 54). Bochner {[}6{]} had proved in 1926 that a function \(f \in \mathcal{C}_b(\mathbf{R})\) is almost periodic if and only if its translates \(x \mapsto f(x - a)\) , for \(a \in \mathbf{R}\) , form a relatively compact subset of \(\mathcal{C}_b(\mathbf{R})\) for the topology of uniform convergence. Bochner and von Neumann observed that if G is a topological group, this characterization visibly provides a notion of almost periodic function (on the right or on the left) on G. In the early 1930s, the theory promised to develop rapidly, in particular under the impetus of von Neumann {[}52{]} and the methods of Peter and Weyl.

Another fertile area appears on the periphery of N. Wiener's work. He introduced, in a famous paper {[}88{]} from 1932, algebraic methods in the study of so-called Tauberian problems—that is, in the study of the asymptotic behavior of a function (or of a sequence) along a filter, given information concerning the behavior of certain weighted means. His approach was motivated by an idea of R. Schmidt.

One of the auxiliary results of Wiener's work ({[}88, Lemma IIe{]}) will particularly strike the minds and inspire numerous developments: if \(f\) is a periodic function whose Fourier series is absolutely convergent, and if \(f\) does not vanish, then the Fourier series of \(1/f\) is absolutely convergent (cf. I, p. 38, example). As early as 1932, R. Paley and N. Wiener {[}89{]} sketched the link between this convergence result, the notion of dual group for a discrete commutative group, and the theory of almost periodic functions.

Against this analytic backdrop, it is algebraic topology that leads L. Pontryagin to give the decisive impetus for the study of characters of locally compact abelian groups. He announces as early as 1932 {[}55{]} the interest of the notion of the character group lies in placing in duality the homology groups of a compact subset of Euclidean space and those of its complement. In a paper {[}58{]} published in 1934, he shows that the character group of a discrete group is compact, and then establishes the duality theorems in this framework. Although Pontryagin uses the results of Peter and Weyl, as well as the Haar measure whose existence had just been established in general, he aims mainly at applications to topology {[}57{]} and is not explicitly concerned with harmonic analysis \(^{(8)}\): the duality theorems are proved by a fine study of the structure of compact abelian groups.

Pontryagin's work immediately aroused great interest. On a suggestion of von Neumann, E. van Kampen extended the duality the following year to all locally compact commutative groups {[}37{]}. Shortly thereafter, he gives a spectacular application of the new invariant analysis to the study of almost periodic functions {[}38{]}, discovered independently by A. Weil during the same period {[}79{]}: If G is a locally compact abelian group, one obtains a compact group \(\widetilde{G}\) beginning by endowing the dual group \(\widehat{G}\) with the discrete topology, and then defining \(\widetilde{G}\) as the compact dual group of the discrete group thus obtained. The group G is then canonically embedded in \(\widetilde{G}\) (cf. II, p. 292, exerc. 54), and a simple application of the Peter--Weyl theory for continuous functions on \(\widetilde{G}\) yields all the results known at the time about almost periodic functions on G.

At the same time, functions of positive type, well known in classical analysis, appear at this period in the study of unitary representations of R. The notion of a universally positive kernel had been studied by J. Mercer {[}45{]} shortly before 1910, in connection with the theory of integral equations, while its analogue for sequences gave rise to a considerable number of works during the same period, related in particular to the moment problem (cf. IV, p. 359, exerc. 21). Meanwhile, functions of positive type on R, already studied systematically by M. Mathias {[}43{]} in 1923, had become, in Bochner's hands, one of the fundamental tools of Fourier analysis and of probability theory {[}7{]}. In 1933, Bochner {[}8{]} and Riesz {[}61{]} independently observed that every diagonal matrix coefficient of a unitary representation of R is a function of positive type on R. This observation enabled them to give a simpler proof of the Stone theorem, and was to have a decisive influence on the development of the general theory (cf. n° 7).

The results of this period were organized by A. Weil in a book completed before 1937 and published in 1940 {[}80{]}. There he sets out in detail the results of Peter--Weyl, Pontryagin, and van Kampen. But in the case of commutative groups, whereas Pontryagin and van Kampen emphasized the structure of groups and duality theorems, Weil systematically develops harmonic analysis, introducing the Fourier transform and proving the Plancherel formula and Bochner's theorem (V, p. 455, th. 5). He insists in particular on the role of convolution (under the name "composition product") and on that of functions of positive type, now defined on an arbitrary group and related to the diagonal matrix coefficients of representations, at least in the case of finite-dimensional representations.

The generality thus achieved in the commutative case would open, a little later, new horizons in number theory. As early as 1936, with a view to introducing algebraic methods into class field theory, C. Chevalley {[}14{]} introduced the group of idèles of a number field, which would later appear as the group of invertible elements of the ring of adèles (cf. AC, VII, p. 221--222, and ÉHM, p. 143). In 1950, J. Tate {[}75{]} would show that harmonic analysis for groups of adèles and idèles makes it possible to recover easily the functional equations of Hecke L-functions, foreshadowing extraordinary developments linking harmonic analysis on adelic groups to the study of automorphic forms.

But one can say that from the end of the 1930s, the principal results of invariant harmonic analysis had been established for compact groups and for locally compact commutative groups. In the commutative case, Weil's proofs still depend (like those of Pontryagin and van Kampen) on a fine knowledge of the structure of the group, that is, roughly speaking, on the classification results of § 2 of II, p. 244. The following decade would show how to dispense with this, thanks to the new methods of the Gelfand school (cf. H. Cartan and R. Godement {[}12{]}).

\histitem{Operator algebras}

We have seen the role that Riesz, in expounding the work of the Hilbert school in 1913, assigned to the algebra \(\mathcal{L}(\mathrm{E})\) of endomorphisms of a Hilbert space, as well as to holomorphic functional calculus and the abstract notion of spectrum. In his 1916 memoir, he already seems to announce normed algebras as the natural framework for spectral theory; after the work of Banach and his school during the decade

of the 1920s, it is hardly surprising that this idea took a central place. The notion of Banach algebra was introduced by M. Nagumo {[}48{]} in 1936, while S. Mazur {[}44{]}, in 1938, showed that the only normed algebra over C which is a field is C itself (I, p. 26, cor. 2).

Toward the middle of the 1930s, in connection with spectral theory, opportunities multiplied for considering algebras of operators on Hilbert spaces abstractly. F. Murray and J. von Neumann, in a series of profound and influential works {[}47{]}, systematically studied the involutive subalgebras of \(\mathcal{L}(\mathrm{E})\) equal to their bicommutant (“von Neumann algebras”). On the other hand, in 1936, S. Steen proposed introducing an axiomatic notion of “algebra of operators” and studying the associated structures in depth {[}70{]}. Moreover, Stone {[}74{]}, motivated by the study of spectral projectors, studied in 1937 the unital algebras all of whose elements are idempotent (“Boolean algebras”: cf. ÉHM, p. 146). He noted that if X is a locally compact topological space, the characteristic functions of open sets and those of complements of everywhere dense sets generate, in \(\mathcal{C}(\mathrm{X})\) , a Boolean algebra A. He then proved that the ideals of A correspond naturally to the closed subsets of the space X, thereby introducing an idea that would prove particularly fruitful.

Beginning in 1939, Gelfand's viewpoint opened new perspectives on these relatively scattered results. His work seems to have been motivated in part by Wiener's ideas, in particular by his theorem on absolutely convergent Fourier series, as well as by the methods of Peter and Weyl. Between 1939 and 1946, in collaboration with M. Naimark, D. Raikov, and G. Shilov, Gelfand established the foundations of the theory of commutative Banach algebras and of the theory of stellar algebras (cf. {[}23{]}, vol. I, p. 169--400). In particular, he gave the classical proof of Wiener's theorem --- generalized by P. Lévy {[}39{]} --- which is found in this book (I, p. 38, example).

In Gelfand's approach, the essential object associated with a commutative Banach algebra A is the space X of maximal ideals of A. Indeed, via the Gelfand--Mazur theorem (I, p. 26, cor. 2), every element of A defines a function on X. It is well known how important this viewpoint would become in the subsequent evolution of algebraic geometry (cf. ÉHM, p. 146--148). This approach was itself motivated by Stone's work on Boolean algebras mentioned above.

It seems that this is L. Loomis, in a course published in 1953 {[}41, p. 53{]}, who first presented Gelfand theory with emphasis on the space of characters of A, defining the Gelfand transform in the sense in which we understand it in this book. The perhaps clearest advantage of this viewpoint is that the topological properties of the space of characters follow immediately from the compactness of the unit ball of the dual of a Banach space in the weak topology.

From his earliest works, Gelfand had constructed the holomorphic functional calculus in one variable in a Banach algebra by means of the Cauchy integral, and had deduced from it that an involutive Banach algebra admits a nontrivial decomposition as a product of rings when the space of its maximal ideals is not connected. The case of an arbitrary Banach algebra was treated only in 1953, when Shilov developed a form of the functional calculus in several variables. {[}68{]}.

As for star algebras, which Gelfand and Naimark had shown as early as 1942 to be realizable as algebras of operators on a Hilbert space (cf. V, p. 442, th. 2), they soon became the object of abundant and profound studies, spurred on by applications to group representations, by the work of Murray and von Neumann, and by the mathematical formalization of quantum mechanics. Let us mention in particular the contributions of I. Segal, who was one of the pioneers of the systematic study of representations of star algebras, in connection with the study of the unitary dual of locally compact groups, which we shall discuss in a few lines.

Although Gelfand, Naimark, Raikov, and Shilov had laid out the essential foundations of the theory of C*-algebras before 1946, certain technical points were improved later. Thus, R. Arens {[}1{]} showed how to avoid invoking the existence of the “Shilov boundary” (cf. I, p. 171, exerc. 26) in order to prove that the Gelfand transform is an involutive morphism. Moreover, the Gelfand–Naimark theory for unital C*-algebras initially added as an axiom the fact that \(1 + x^{*}x\) is invertible for all \(x\). They conjectured, however, that this axiom was unnecessary; the proof of this assertion, elementary but delicate, was only provided around 1952, following the work of M. Fukamiya {[}21{]} taken up by I. Kaplansky {[}62{]}.

\histitem{Representations of locally compact groups}

Until 1939, the study of representations of topological groups was limited to the compact case and the commutative case. Physicists, for their part, had found in quantum mechanics many reasons to study irreducible unitary representations in order to seek Hilbert spaces capable of modeling the states of elementary particles. P. Dirac, building on work by E. Majorana, had pointed out around 1936 the interest, in this perspective, of the Lorentz group SO(3,1) and the Poincaré group SO(3,1) × \(R^{4}\). Their methods remained far removed from those of the mathematicians of the time. But in 1939, E. Wigner {[}90{]} opened a breach when he classified the irreducible unitary representations of the Poincaré group, building on the work of Murray and von Neumann.

It is to Gelfand that we owe, it seems, having glimpsed a general path toward the study of representations of locally compact groups. In a memoir written in 1942 (cf. {[}23{]}, vol. II, p. 3--17), he notes with Raikov that such a theory is made possible by the link between unitary representations and functions of positive type. They prove the existence of a Hilbertian realization for every function of positive type (V, p. 432, th. 1) and deduce from it that irreducible unitary representations separate the points of G (V, p. 454, th. 4). The positive linear forms on the involutive algebra L \(^{1}\) (G) play an essential role: Gelfand and Raikov show the link that connects them to functions of positive type (cf. V, p. 448, prop. 13). The enveloping C*-algebra of an involutive algebra is presented (without a name) in {[}24, § 48{]}, although it is not yet directly a question of associating with G the algebra Stell(G) (cf. 9, def. of I, p. 125).

After 1945, the study of unitary representations developed by leaps and bounds. The interaction with C*-algebras, as well as with the analysis of the foundations of quantum theory, quickly provided valuable general results: around 1947, Segal associated with a locally compact group various \(^{(9)}\) stellar algebras {[}67{]} and wove the link between their representations and those of the group under consideration. Shortly thereafter, G. Mackey extracts the notion of induced representation {[}42{]} which will be omnipresent in the concrete results on unitary representations.

At the same time, the theory is considerably stimulated by the study of examples that reveal its variety and depth. Gelfand and Naimark, studying the example of SL(2, C), isolated in 1947 a special case of the notion of induced representation, then showed that it is the key to the study of the representations of this group {[}23, vol. II, p. 41--124{]}. In the same year, V. Bargmann {[}5{]} studied the group SL(2, R) and discovered the existence of a countable family of square-integrable representations, called the "discrete series," as well as the orthogonality relations between their matrix coefficients (cf. prop. 8 of V, p. 424). As Godement immediately realized {[}25{]}, the Bargmann orthogonality relations are satisfied by all square-integrable representations of locally compact groups.

Soon the study of particular classes of groups would again, and for a long time, take precedence over abstract theory. Harish-Chandra, in particular, would undertake the general study of representations of reductive Lie groups --- an epic task that we cannot describe here, any more than the prodigious consequences that Langlands would foresee for number theory.

\specialchapter{BIBLIOGRAPHY}

{[}1{]} R. ARENS - "On a theorem of Gelfand and Neumark", Proc. Natl. Acad. Sci. USA 32 (1946), p. 237-239, Zbl 0060.27101.

{[}2{]} M. F. ATIYAH and I. M. SINGER -- "The index of elliptic operators on compact manifolds", Bull. Amer. Math. Soc. 69 (1963), p. 422--433, Zbl 0118.31203; Michael Atiyah Collected Works, vol. 3, Clarendon Press, 1988, p. 13--24.

{[}3{]} F. V. ATKINSON -- "Normal solvability of linear equations in normed spaces", Mat. Sb., Nov. Ser. 28 (1951), p. 3--14, Zbl 0042.12001 (in Russian).

{[}4{]} S. BANACH - Theory of linear operations, Warsaw, 1932, Zbl 0005.20901.

{[}5{]} V. BARGMANN - "Irreducible unitary representations of the Lorentz group", Ann. of Math. (2) 48 (1947), p. 568-640, Zbl 0045.38801.

{[}6{]} S. BOCHNER -- "Contributions to the theory of almost periodic functions. II: Functions of several variables", Math. Ann. 96 (1926), p. 383--409, JfM 52.0265.01; Collected papers, vol. 2, Amer. Math. Soc., 1992, p. 79--106.

{[}7{]} \_\_\_\_ Lectures on Fourier integrals, Leipzig, 1932, Zbl 0006.11001.

{[}8{]} \_\_\_\_ "Spectral decomposition of linear families of unitary operators", Sitzungsber. Preuß. Akad. Wiss., Phys.-Math. Kl. (1933), p. 371--376, Zbl 0007.16603; Collected papers, vol. 1, Amer. Math. Soc., 1992, p. 579--586.

{[}9{]} H. BOHR - "On the theory of almost periodic functions. I. A generalization of the theory of Fourier series", Acta Math. 45 (1924), p. 29-127, JfM 50.0196.01.

{[}10{]} J. W. CALKIN - "Two-sided ideals and congruences in the ring of bounded operators in Hilbert space", Ann. of Math. (2) 42 (1941), p. 839-873, Zbl 0063.00692.

{[}11{]} T. CARLEMAN - On singular integral equations with real and symmetric kernel, Uppsala, 1923, JfM 49.0272.01.

{[}12{]} H. CARTAN and R. GODEMENT -- "Duality theory and harmonic analysis in locally compact abelian groups", Ann. Sci. Éc. Norm. Supér. (3) 64 (1947), p. 79--99, Zbl 0033.18801; Henri Cartan's Works, vol. III, Springer, 1979, p. 1203--1223.

{[}13{]} H. CARTAN and J-P. SERRE -- "A finiteness theorem concerning compact analytic manifolds", C. R. Acad. Sci., Paris 237 (1953), p. 128--130, Zbl 0050.17701; Henri Cartan's Works, vol. II, Springer, 1979, p. 684--686.

{[}14{]} C. CHEVALLEY - "Generalization of class field theory for infinite extensions", J. Math. Pures Appl. (9) 15 (1936), p. 359-371, Zbl 0015.15101.

{[}15{]} J. DIEUDONNÉ -- "On homomorphisms of normed spaces", Bull. Sci. Math. (2) 67 (1943), p. 72--84, Zbl 0028.23301; Selected mathematical works, Hermann, 1981, p. 268--281.

{[}16{]} J. DIEUDONNÉ - History of functional analysis, North-Holland Publ. Co., 1981, Zbl 0478.46001.

{[}17{]} P. L. DIRICHLET -- "Proof of the theorem that every unbounded arithmetic progression, whose first term and difference are integers without common factor, contains infinitely many primes", Abh. Preuß. Akad. Wiss. 8 (1837), p. 45--81; Works, vol. 1, Reimer (Berlin), 1889, p. 315--356.

{[}18{]} P. ENFLO - "A counterexample to the approximation problem in Banach spaces", Acta Math. 130 (1973), p. 309-317, Zbl 0267.46012.

{[}19{]} \_\_\_\_ "On the invariant subspace problem for Banach spaces", Acta Math. 158 (1987), no. 3-4, p. 213--313, Zbl 0663.47003.

{[}20{]} J. M. G. FELL - "The dual spaces of C*-algebras", Trans. Amer. Math. Soc. 94 (1960), p. 365-403, Zbl 0090.32803.

{[}21{]} M. FUKAMIYA - "On a theorem of Gelfand and Neumark and the B*-algebra", Kumamoto J. Sci., Math. 1 (1952), p. 17-22, Zbl 0049.08605.

{[}22{]} C. F. GAUSS - Disquisitiones Arithmeticae, Göttingen, 1801.

{[}23{]} I. M. GELFAND - Collected papers, 3 vol., Springer, 1987.

{[}24{]} I. M. GELFAND and M. A. NAIMARK -- "Normed rings with involution and their representations", Izv. Akad. Nauk SSSR, Ser. Mat. 12 (1948), p. 445--480, Zbl 0031.03403; I. M. Gelfand Collected papers, vol. I, Springer, 1987, p. 364--400 (in Russian).

{[}25{]} R. GODEMENT -- "On the orthogonality relations of V. Bargmann. II. General proof", C. R. Acad. Sci. Paris 225 (1947), p. 657--659, Zbl 0029.19907.

{[}26{]} I. C. GOHBERG -- "On linear equations in normed spaces", Dokl. Akad. Nauk SSSR, n. Ser. 76 (1951), p. 477--480, Zbl 0042.11903 (in Russian).

{[}27{]} P. R. HALMOS - "What does the spectral theorem say?", Amer. Math. Monthly 70 (1963), p. 241-247, Zbl 0132.35606; Selecta, vol. II, Springer, 1983, p. 59-66.

{[}28{]} W. HEISENBERG -- "On a modification of the formal rules of quantum theory in the problem of the anomalous Zeeman effects", Z. Phys. 26 (1924), p. 291--307, JfM 51.0740.07; Collected Works: Scientific original papers, vol. 1, Springer, 1985, p. 306--325.

{[}29{]} E. HELLINGER and O. TOEPLITZ -- "Integral equations and equations with infinitely many unknowns", Enzyklopädie der mathematischen Wissenschaften, art. II C 13, Leipzig (Teubner), 1927, JfM 53.0350.01.

{[}30{]} E. HILB - "On integral representation of arbitrary functions", Math. Ann. 66 (1909), p. 1-66, JfM 39.0380.03.

{[}31{]} D. HILBERT - "Foundations of a general theory of linear integral equations. First communication.", Nachr. Ges. Wiss. Göttingen, Math.-Phys. Kl. (1904), p. 49-91, JfM 35.0378.02.

{[}32{]} \_\_\_\_ "Foundations of a general theory of linear integral equations. Fourth communication.", Nachr. Ges. Wiss. Göttingen, Math.-Phys. Kl. (1906), p. 157--227, JfM 37.0351.03.

{[}33{]} T. H. HILDEBRANDT - "On completely continuous linear transformations", Acta Math. 51 (1928), p. 311-318, JfM 54.0427.03.

{[}34{]} E. HILLE - Functional analysis and semi-groups, Colloq. Publ., vol. 31, Amer. Math. Soc., 1948, Zbl 0033.06501.

{[}35{]} M. JAMMER - The conceptual development of quantum mechanics, McGraw-Hill, 1966.

{[}36{]} M. KAC - "Can one hear the shape of a drum?", Am. Math. Mon. 73 (1966), p. 1-23, Zbl 0139.05603.

{[}37{]} E. R. VAN KAMPEN - "Locally bicompact abelian groups and their character groups", Ann. of Math. (2) 36 (1935), p. 448-463, Zbl 0011.39401.

{[}38{]} \_\_\_\_ "Almost periodic functions and compact groups", Ann. of Math. (2) 37 (1936), p. 78--91, Zbl 0014.01702.

{[}39{]} P. LÉVY - "On the absolute convergence of Fourier series", Compos. Math. 1 (1934), p. 1-14, Zbl 0008.31201.

{[}40{]} V. I. LOMONOSOV - "Invariant subspaces for the family of operators which commute with a completely continuous operator", Funct. Anal. Appl. 7 (1974), p. 213-214, Zbl 0293.47003.

{[}41{]} L. H. LOOMIS - An introduction to abstract harmonic analysis, Van Nostrand Co., 1953, Zbl 0052.11701.

{[}42{]} G. W. MACKEY - "Imprimitivity for representations of locally compact groups. I", Proc. Nat. Acad. Sci. U.S.A. 35 (1949), p. 537-545, Zbl 0035.06901.

{[}43{]} M. MATHIAS - "On positive Fourier integrals", Math. Z. 16 (1923), p. 103--125, JfM 49.0207.03.

{[}44{]} S. MAZUR - "On linear rings", C. R. Acad. Sci., Paris 207 (1938), p. 1025-1027, Zbl 0020.20101.

{[}45{]} J. MERCER - "Functions of positive and negative type, and their connection with the theory of integral equations", Philos. Trans. R. Soc. Lond., Ser. A 209 (1909), p. 415-446, JfM 40.0408.02.

{[}46{]} S. G. MIKHLIN - "On the resolution of linear equations in Hilbert spaces", Dokl. Akad. Nauk SSSR (N. S.) 57 (1947), p. 11-12, Zbl 0029.05101 (in Russian).

{[}47{]} F. J. MURRAY and J. VON NEUMANN - "On rings of operators", Ann. of Math. (2) 37 (1936), p. 116-229, Zbl 0014.16101; John von Neumann Collected works, vol. III, Pergamon Press, 1961, p. 6-120.

{[}48{]} M. NAGUMO - "Some analytic investigations in linear metric rings", Jpn. J. Math. 13 (1936), p. 61-80, Zbl 3024132.

{[}49{]} J. VON NEUMANN -- "Mathematical foundations of quantum mechanics", Göttingen Nachrichten 1927 (1927), p. 1--57, JfM 53.0848.03; Collected works, vol. I, Pergamon Press, 1961, p. 151--207.

{[}50{]} \_\_\_\_ "General eigenvalue theory of Hermitian functional operators", Math. Ann. 102 (1929), p. 49--131, JfM 55.0824.02; Collected works, vol. II, Pergamon Press, 1961, p. 3--85.

{[}51{]} \_\_\_\_ "On the algebra of functional operations and the theory of normal operators", Math. Ann. 102 (1929), p. 370--427, JfM 55.0825.02; Collected works, vol. II, Pergamon Press, 1961, p. 86--143.

{[}52{]} \_\_\_\_ "Almost periodic functions in a group. I", Trans. Amer. Math. Soc. 36 (1934), p. 445--492, Zbl 0009.34902; Collected works, vol. II, Pergamon Press, 1961, p. 454--501.

{[}53{]} F. NOETHER -- "On a class of singular integral equations", Math. Ann. 82 (1920), p. 42--63, JfM 47.0369.02.

{[}54{]} F. PETER and H. WEYL -- "The completeness of the primitive representations of a closed continuous group", Math. Ann. 97 (1927), p. 737--755, JfM 53.0387.02; Hermann Weyl Gesammelte Abhandlungen, Bd. III, Springer, 1968, p. 58--75.

{[}55{]} L. PONTRYAGIN -- "The general duality theorem for closed sets", Verhandlungen Kongreß Zürich, 2, p. 195-197, 1932, JfM 58.0646.08.

{[}56{]} \_\_\_\_ "Almost periodic functions and analysis situs", C. R. Acad. Sci., Paris 196 (1933), p. 1201--1203, Zbl 0006.42801.

{[}57{]} \_\_\_\_ "The general topological theorem of duality for closed sets", Ann. of Math. (2) 35 (1934), no. 4, p. 904--914, Zbl 0010.18004.

{[}58{]} \_\_\_\_ "The theory of topological commutative groups", Ann. of Math. (2) 35 (1934), p. 361--388, Zbl 0009.15601.

{[}59{]} F. RIESZ - Les systèmes d'équations linéaires à une infinité d'inconnues, Gauthier-Villars, 1913, JfM 44.0401.01.

{[}60{]} \_\_\_\_ "On linear functional equations", Acta Math. 41 (1916), p. 71--98, JfM 46.0635.01 ; Œuvres complètes, t. II, Gauthier-Villars, 1960, p. 1053--1080.

{[}61{]} \_\_\_\_ "On theorems of Stone and Bochner", Acta Litt. Sci. Szeged 6 (1933), p. 184--198, Zbl 0007.30903; Œuvres complètes, t. II, Gauthier-Villars, 1960, p. 1135--1149.

{[}62{]} J. SCHATZ - "Review of the article On a theorem of Gelfand and Neumark and the B*-algebra by M. Fukamiya", Math. Reviews 14 (1952), p. 884.

{[}63{]} J. SCHAUDER - "On linear, completely continuous functional operations", Studia Math. 2 (1930), p. 183-196, JfM 56.0354.01.

{[}64{]} A. SCHUSTER -- "The genesis of spectra", Report on the fifty-second meeting of the British association for the advancement of Science, J. Murray (London), 1882, p. 120--143.

{[}65{]} L. SCHWARTZ - "Generalization of the notion of function, of derivation, of Fourier transformation and mathematical and physical applications", Ann. Univ. Grenoble, Sect. Sci. Math. Phys., n. Ser. 21 (1946), p. 57-74, Zbl 0060.27504; Œuvres scientifiques, t. I, Soc. Math. France, 2011, p. 61-88.

{[}66{]} \_\_\_\_ "Homomorphisms and completely continuous mappings", C. R. Acad. Sci., Paris 236 (1953), p. 2472--2473, Zbl 0050.33301 ; Œuvres scientifiques, t. I, Soc. Math. France, 2011, p. 345--346.

{[}67{]} I. E. SEGAL - "The group algebra of a locally compact group", Trans. Amer. Math. Soc. 61 (1947), p. 69-105, Zbl 0032.02901.

{[}68{]} G. E. SHILOV -- "On the decomposition of a commutative normed ring into a direct sum of ideals", Mat. Sb., Nov. Ser. 32 (1953), p. 353--364, Zbl 0052.34203 (in Russian).

{[}69{]} S. SOBOLEV - "New method to solve the Cauchy problem for normal hyperbolic linear equations", Rec. Math. Moscou, n. Ser. 1 (1936), p. 39-71, Zbl 0014.05902.

{[}70{]} S. W. P. STEEN - "An introduction to the theory of operators : I. Real operators. Modulus", Proc. Lond. Math. Soc. (2) 41 (1936), p. 361-392, Zbl 0014.16201.

{[}71{]} T.-J. STIELTJES - "Researches on continued fractions", Ann. Fac. Sci. Toulouse Sci. Math. Sci. Phys. 8 (1894), no. 4, p. J1-J122, JfM 25.0326.01.

{[}72{]} M. H. STONE - Linear transformations in Hilbert space and their applications to analysis, Colloq. Publ., Am. Math. Soc., vol. 15, Amer. Math. Soc., 1932, Zbl 0005.40003.

{[}73{]} \_\_\_\_ "On one-parameter unitary groups in Hilbert space", Ann. of Math. (2) 33 (1932), p. 643--648, Zbl 0005.16403.

{[}74{]} \_\_\_\_ "Applications of the theory of Boolean rings to general topology", Trans. Am. Math. Soc. 41 (1937), p. 375--481, Zbl 0017.13502.

{[}75{]} J. T. TATE — “Fourier analysis in number fields, and Hecke's zetafunctions”, Thesis, Princeton, 1950; Collected works, part I, Amer. Math. Soc., 2016, p. 1–43, Zbl 1407.01030.

{[}76{]} O. TOEPLITZ — “Das algebraische Analogon zu einem Satze von Fejér”, Math. Z. 2 (1918), p. 187–197, JfM 46.0157.02.

{[}77{]} H. WEBER — “Beweis des Satzes, dass jede eigentlich primitive quadratische Form unendlich viele Primzahlen darzustellen fähig ist.”, Math. Ann. 20 (1882), p. 301–330, JfM 14.0141.01.

{[}78{]} \_\_\_\_ “Theorie der Abel'schen Zahlkörper”, Acta Math. 9 (1887), p. 105–130, JfM 18.0055.04.

{[}79{]} A. WEIL — “Sur les fonctions presque périodiques”, C. R. Acad. Sci., Paris 200 (1935), p. 38–40, Zbl 0010.35303; Œuvres scientifiques, vol. I, Springer, 1979, p. 101–103.

{[}80{]} \_\_\_\_ L'intégration dans les groupes topologiques et ses applications, Actualités scientifiques et industrielles, vol. 869, Hermann, 1940, Zbl 0063.08195.

{[}81{]} H. WEYL — “Über beschränkte quadratische Formen, deren Differenz vollstetig ist”, Rend. Circ. Mat. Palermo 27 (1909), p. 373–392, JfM 40.0395.01; Gesammelte Abhandlungen, Bd. I, Springer, 1968, p. 175–194.

{[}82{]} \_\_\_\_ “Über gewöhnliche Differentialgleichungen mit Singularitäten und die zugehörigen Entwicklungen willkürlicher Funktionen”, Math. Ann. 68 (1910), p. 220–269, JfM 41.0343.01; Gesammelte Abhandlungen, Bd. I, Springer, 1968, p. 248–297.

{[}83{]} \_\_\_\_ “Theorie der Darstellung kontinuierlicher halbeinfacher Gruppen durch lineare Transformationen. I, II, III, Nachtrag”, Math. Z. 23 (1925), p. 271–309, JfM 51.0319.01; Gesammelte Abhandlungen, Bd. II, Springer, 1968, p. 543–647.

{[}84{]} \_\_\_\_ “Integralgleichungen und fastperiodische Funktionen”, Math. Ann. 97 (1926), p. 338–356, JfM 52.0260.02; Gesammelte Abhandlungen, Bd. III, Springer, 1968, p. 38–57.

{[}85{]} \_\_\_\_ “Quantenmechanik und Gruppentheorie”, Z. Phys. 46 (1927), p. 1–46, JfM 53.0848.02; Gesammelte Abhandlungen, Bd. III, Springer, 1968, p. 90–135.

{[}86{]} \_\_\_\_ Gruppentheorie und Quantenmechanik, S. Hirzel (Leipzig), 1928, Zbl 0537.22023.

{[}87{]} E. WEYR — “Zur Theorie der bilinearen Formen”, Monatsh. Math. Phys. 1 (1890), p. 163–236, JfM 22.0141.02.

{[}88{]} N. WIENER — “Tauberian theorems”, Ann. of Math. (2) 33 (1932), p. 1–100, Zbl 0004.05905; Selected Papers, MIT Press, 1964, p. 261–360.

{[}89{]} N. WIENER and R. E. A. C. PALEY — “Analytic properties of the characters of infinite abelian groups”, Verhandlungen Kongreß Zürich, 2, p. 95, 1932, JfM 58.0445.02.

{[}90{]} E. WIGNER — “On unitary representations of the inhomogeneous Lorentz group”, Ann. of Math. (2) 40 (1939), p. 149–204, Zbl 0020.29601; Collected works, part A, Springer, 1993, p. 334–389.

{[}91{]} W. WIRTINGER — “Beiträge zu Riemann's Integrationsmethode für hyperbolische Differentialgleichungen, und deren Anwendungen auf Schwingungsprobleme”, Math. Ann. 48 (1897), p. 365–389, JfM 27.0281.01.

{[}92{]} S. L. WORONOWICZ — “Unbounded elements affiliated with C*-algebras and noncompact quantum groups”, Comm. Math. Phys. 136 (1991), no. 2, p. 399–432, Zbl 0743.46080.

{[}93{]} B. YOOD — “Properties of linear transformations preserved under addition of a completely continuous transformation”, Duke Math. J. 18 (1951), p. 599–612, Zbl 0043.11901.

\specialchapter{INDEX OF NOTATIONS}

\(\mathcal{L}^{\mathrm{c}}(\mathrm{E};\mathrm{F})\) (compact linear maps) 2 \(\mathrm{E}_{(\mathbf{C})}\) (complexified vector space) 2 \(u_{(\mathbf{C})}\) (complexified linear map) 3 \(\mathcal{L}^{\mathrm{f}}(\mathrm{E};\mathrm{F})\) (finite-rank linear maps) 4 \(\mathcal{C}_0(\mathrm{X};\mathrm{K}), \mathcal{C}_0(\mathrm{X})\) (continuous functions tending to 0 at infinity) 20 \(\mathcal{N}^{p,q}(\mathrm{X} \times \mathrm{Y})\) (kernel space) 27 \(u_k\) (integral operator with kernel) 29 \(k_\circ\) (map \(y \mapsto k_y\) where \(k_y(x) = k(x,y)\) ) 30 \(\mathcal{L}^{p',q'}(\mathrm{X},\mathrm{Y},\mu,\nu)\) (kernel space) 31 \(u \equiv v\) (congruence modulo \(\mathcal{L}^{\mathrm{f}}(\mathrm{E};\mathrm{F})\) ) 40 \(\mathcal{F}(\mathrm{E};\mathrm{F})\) (Fredholm maps) 41\\
ind(u) (index of a Fredholm map) 43 \(\mathcal{I}(\mathrm{E};\mathrm{F})\) (isomorphisms of E onto F) 57 \(\mathcal{M}\mathcal{D}(\mathrm{E};\mathrm{F})\) (injective direct morphisms) 57 \(\mathcal{E}\mathcal{D}(\mathrm{E};\mathrm{F})\) (surjective direct morphisms) 57 \(\mathcal{R}(\mathrm{E})\) (Riesz endomorphisms) 60\\
((u)) (conorm of u) 61 \(\mathcal{M}(\mathrm{E};\mathrm{F})\) (strict injective morphisms) 67 \(\mathcal{Q}\mathcal{M}(\mathrm{E};\mathrm{F})\) (strict morphisms with finite-dimensional kernel) 67 \(\mathcal{E}(\mathrm{E};\mathrm{F})\) (continuous surjective morphisms) 70 \(\mathcal{Q}\mathcal{E}(\mathrm{E};\mathrm{F})\) (morphisms whose image has finite codimension) 70 \(\mathcal{C}alk(\mathrm{E})\) (Calkin algebra of a Banach space E) 75 \(\mathrm{N}_{\lambda}\) (nilspace) 78 \(\mathrm{I}_{\lambda}\) (conilspace) 78 \(e_{\lambda}(u)\) (spectral projector) 82 \(m_{\lambda}(u)\) (spectral multiplicity) 82 \(\mathrm{Sp}_{\mathrm{s}}(u)\) (sensitive points of the spectrum) 83 \(\mathrm{Sp}_{\mathrm{e}}(u)\) (essential spectrum) 83 \(\mathrm{Sp}_{n}(u)\) (elements of the spectrum such that \(u - \lambda 1_{\mathrm{E}}\) has index n) 86 \(\mathrm{Sp}_{\omega}(u)\) (complement of the union of the \(\mathrm{Sp}_{n}(u)\) ) 86 \(\mathcal{D}_{\mathrm{B}}(\mathrm{E})\) (endomorphisms diagonal in the basis B) 146 \(\mathrm{I}_{\mathrm{E}}\) (dimensions of finite-dimensional subspaces F ≠ E) 151 \(\lambda_n(u)\) (decreasing sequence of eigenvalues of u) 152

\(r_{\mathrm{F}}(u), \mathrm{R}_{\mathrm{F}}(u)\) (Rayleigh quotients) \ldots.. 153 \(\mathcal{F}_n\) (subspaces of dimension \(n\) ) \ldots.. 153 \((\alpha_n(u))_{n \in \mathrm{J}}\) (sequence of singular values) \ldots.. 156 \((\alpha_n(u))_{n \in \mathrm{I_E}}\) (extended sequence of singular values) \ldots.. 156 \(\mathcal{L}_1(\mathrm{E})\) (endomorphisms of finite trace) \ldots.. 164 \(\|u\|_1\) (norm on \(\mathcal{L}_1(\mathrm{E};\mathrm{F})\) ) \ldots.. 167 \(\mathcal{L}_1(\mathrm{E};\mathrm{F})\) (nuclear linear maps) \ldots.. 169 \(\mathcal{L}_{\mathrm{u}}(\mathrm{X})\) (universally measurable functions) \ldots.. 183 \(\mathcal{L}_{\mathrm{u}}^{\infty}(\mathrm{X})\) (bounded universally measurable functions) \ldots.. 183 \(m_g, \widetilde{m}_g\) (multiplication endomorphisms) \ldots.. 186 \(\mathcal{G}_{\mathrm{A}}\) (Gelfand transformation) \ldots.. 190 \(\mathcal{H}_{\mathrm{A}}\) (inverse of the Gelfand transformation) \ldots.. 190 \(\mathcal{F}\) (Fourier transformation) \ldots.. 196 \(\overline{\mathcal{F}}\) (Fourier cotransformation) \ldots.. 196 \(\mathcal{D}(\mathrm{U})\) (test functions on U) \ldots.. 200 \(\mathcal{D}_{\mathbf{R}}(\mathrm{U})\) (real-valued test functions) \ldots.. 201 \(\mathcal{D}'(\mathrm{U})\) (distributions on U) \ldots.. 203 \(\partial^\alpha f\) (partial derivative of a distribution) \ldots.. 208 \(\mathcal{S}(\mathbf{R}^n)\) (Schwartz functions on \(\mathbf{R}^n\) ) \ldots.. 209 \(\mathcal{S}'(\mathbf{R}^n)\) (tempered distributions on \(\mathbf{R}^n\) ) \ldots.. 215 \(\mathcal{M}^t(\mathbf{R}^n)\) (space of tempered measures) \ldots.. 216 \(W^{k,p}(U)\) (Sobolev space of index \(k\) and exponent \(p\) ) \ldots.. 222 \(W_0^{k,p}(U)\) (closure of \(\mathcal{D}(U)\) in \(W^{k,p}(U)\) ) \ldots.. 222 \(H^k(U)\) (Sobolev space of index \(k\) and exponent 2) \ldots.. 222 \(H_0^k(U)\) (closure of \(\mathcal{D}(U)\) in \(H^k(U)\) ) \ldots.. 222 \(\mathcal{P}(\mathrm{E};\mathrm{F})\) (partial operators from E to F) \ldots.. 224 \(\mathcal{P}(\mathrm{E})\) (partial operators on E) \ldots.. 224 dom( \(u\) ) (domain of \(u\) ) \ldots.. 224 \(1_{\mathrm{D}}\) (identical partial operator with domain D) \ldots.. 225 \(u \circ v\) (composition of partial operators \(u\) and \(v\) ) \ldots.. 225 \(u \subset v\) ( \(v\) is an extension of \(u\) ) \ldots.. 226 \((u_i)_{i \in I}\) (product partial operator) \ldots.. 226 \(u + v\) (sums of partial operators \(u\) and \(v\) ) \ldots.. 226 \(\overline{u}\) (closure of \(u\) ) \ldots.. 228 \(E_u\) (normed space defined by the domain of \(u\) ) \ldots.. 230 \(\|x\|_u\) (norm on the space \(E_u\) ) \ldots.. 230 \((x | y)_u\) (inner product defining the Hilbert space \(E_u\) ) \ldots.. 230 \(m_g, \widetilde{m}_g\) (multiplication partial operator) \ldots.. 231 \(u^*\) (adjoint of \(u\) ) \ldots.. 235 \(\mathcal{A}(\mathrm{E})\) (self-adjoint partial operators) \ldots.. 238 \(^t\mathrm{D}\) (transposed differential operator) \ldots.. 242 \(H_D\) (domain of the differential operator \(D_+\) ) \ldots.. 242 \(D_-\) (differential operator with domain \(\mathcal{D}(U)\) ) \ldots.. 242 \(D_+\) (differential operator with domain \(H_D\) ) \ldots.. 242 \(R(u,\lambda)\) (resolvent of \(u\) ) \ldots.. 244 \(\mathrm{PSp}_{\varepsilon}(u)\) (\(\varepsilon\)-pseudo-spectrum) \ldots.. 250 \(b(u)\) (bornification of \(u\) ) \ldots.. 262 \(\Omega(\mathrm{E})\) (set of \(v\) such that \(1_{\mathrm{E}} - v^*v\) is positive and injective) \ldots.. 264 \(B(v)\) (partial operator \(v\circ((1_{\mathrm{E}} - v^*v)^{1/2})^{-1}\) ) \ldots.. 264

\(\beta (z) = z / \sqrt{1 + |z|^2}\) 266 \(D_f\) (domain of \(f(u)\) ) 271 \(f(u)\) (universally measurable functional calculus) 272 \(p_A\) (spectral projector defined by A) 277 \(u^\alpha\) (functional calculus) 289 \(|u|\) (absolute value of \(u\) ) 290 \(q_u\) (positive partial form associated with \(u\) ) 291 \(E_q\) (pre-Hilbert space) 292 \((x|y)_q\) (inner product on \(E_q\) ) 292 \(\| x \| _q\) (norm on \(E_q\) ) 292 \(u_q\) (partial operator representing \(q\) ) 292 \(\mathrm{Sp}_{\mathrm{s}}(u)\) (sensitive spectrum of \(u\) ) 297 \(\mathrm{Sp}_{\mathrm{e}}(u)\) (essential spectrum of \(u\) ) 297 \(\mathrm{Sp_h}(u)\) (upper part of the spectrum) 298 \(\varrho_{e}\) (lower bound of the essential spectrum) 298 \(\mathrm{Sp_b}(u)\) (lower part of the spectrum) 298 \(\mathcal{F}_n\) (subspaces of dimension \(n\) ) 299 \(\mathcal{F}_n^u\) (subspaces of \(\mathrm{dom}(u)\) of dimension \(n\) ) 299 \(\widetilde{r}_{\mathrm{F}}(u), \widetilde{\mathrm{R}}_{\mathrm{F}}(u)\) (Rayleigh quotient) 300 \(\|v\|_u\) (relative norm of \(v\) with respect to \(u\) ) 306\\
vp (Cauchy principal value) 337 \(\chi_\varrho\) (character of a representation) 374 \(\mathrm{Res}_{H}^{G}(\varrho)\) (restriction of \(\varrho\) to H) 375 \(\mathrm{Hom}_{G}(\varrho_1,\varrho_2)\) (spaces of G-morphisms) 375 \(\varrho^n = n \varrho\) (sum of copies of the representation \(\varrho\) ) 376 \(\check{\varrho}\) (contragredient representation) 377 \(\mathrm{Sesq_G}(\varrho, \pi), \mathrm{Sesq_G(E,F)}\) (G-invariant sesquilinear forms) 380 \(\overline{\varrho}\) (conjugate representation) 382 \(\varrho_1 \boxtimes \varrho_2\) (external tensor product) 385 \(\varrho_1 \otimes \varrho_2\) (tensor product) 385 \(\varrho^{\otimes n}\) (tensor power of \(\varrho\) ) 385 \(\Upsilon(G)\) (matrix coefficients) 386 \(\Theta(G)\) (finite-dimensional matrix coefficients) 386 \(\operatorname{cl}(\pi)\) (class of a unitary representation) 392 \(\widehat{G}\) (unitary dual) 393 \(M_\pi (\varrho)\) (isotypic component) 394 \(\mathcal{M}_c(G)\) (measures with compact support) 400 \(\mathcal{M}^1(G)\) (bounded measures) 401 \(\gamma_{G}^{(p)}\) (left regular representation) 405 \(\delta_{G}^{(p)}\) (right regular representation) 405 \(\varrho_{G}^{(p)}\) (biregular representation) 406 \(\gamma_{G}\) (left regular representation in \(L^2\) ) 407 \(\delta_{G}\) (right regular representation in \(L^2\) ) 407 \(\varrho_{G}\) (biregular representation in \(L^2\) ) 407 \(\mathcal{F}_{\pi}(G)\) (π-equivariant functions) 408 \(\mathcal{K}_{\pi}(G)\) (continuous functions with compact support modulo H in) \(\mathcal{F}_{\pi}(G)\) ) \ldots{} 408 \(N_p(f)\) (semi-norm on) \(\mathcal{F}_π(G)\) ) \ldots{} 408 \(\mathcal{L}_{\pi}^{p}(G)\) (π-equivariant functions of power \(p^e\) integrable modulo H) \ldots{} 409

\(L_{\pi}^{p}(G)\) (Banach space associated with \(\mathcal{L}_{\pi}^{p}(G)\) ). \ldots.. 410 \(\mathcal{L}_{\pi}^{2}(G)\) (\(\pi\)-equivariant functions square-integrable modulo H) \ldots.. 415 \(L_{\pi}^{2}(G)\) (Hilbert space associated with \(\mathcal{L}_{\pi}^{2}(G)\) ). \ldots.. 415 \(\operatorname{Ind}_{H}^{G}(\pi)\) (representation of G induced by the representation π of H). \ldots.. 416 \(\gamma_{G,\chi}\) (left regular representation on \(\mathcal{F}_{\chi}(G)\) or \(L_{\chi}^{2}(G)\) ). \ldots.. 417 \(\delta_{G,\chi}\) (right regular representation on \(\mathcal{F}_{\chi}(G)\) or \(L_{\chi}^{2}(G)\) ). \ldots.. 417 \(\gamma_{G,\chi}^{(p)}\) (left regular representation in \(L_{\chi}^{p}(G)\) ). \ldots.. 417 \(\delta_{G,\chi}^{(p)}\) (right regular representation in \(L_{\chi}^{p}(G)\) ). \ldots.. 417 \(\gamma_{G,\chi}\) (unitary left regular representation in \(L_{\chi}^{2}(G)\) ). \ldots.. 417 \(\delta_{G,\chi}\) (unitary right regular representation in \(L_{\chi}^{2}(G)\) ). \ldots.. 417 \(\varrho_{G,\chi}\) (unitary biregular representation in \(L_{\chi}^{2}(G)\) ). \ldots.. 417 \(c_{Z}(\pi)\) (formal degree). \ldots.. 423 \(\mathrm{Noy}_{+}(\mathrm{X})\) (universally positive kernels). \ldots.. 434 \(A_{+}^{\prime}\) (continuous positive linear forms). \ldots.. 436 \(\mathrm{Pos}(G)\) (functions of positive type). \ldots.. 443 \(\mathrm{Pos}_{1}(G)\) (functions of positive type such that \(\varphi(e)=1\) ). \ldots.. 443 \(\varrho_{G}(\pi)\) (\((\overline{\pi}\boxtimes\pi)\)-isotypic component of \(\varrho_{G}\) ). \ldots.. 462 \(\Theta(G^{\sharp})\), \(Z\Theta(G)\) (central matrix coefficients). \ldots.. 467 \(R(G)\) (Grothendieck ring). \ldots.. 469 \(\mathrm{F}(\widehat{G})\) (product algebra of the End(Eπ)). \ldots.. 470 \(\mathrm{F}_{b}(\widehat{G})\) (subalgebra of bounded families). \ldots.. 470 \(\|u\|_{\mathrm{F}_{b}}\) (norm in \(\mathrm{F}_{b}(\widehat{G})\)). \ldots.. 470 \(\mathrm{F}_{0}(\widehat{G})\) (subalgebra of families tending to 0). \ldots.. 470 \(\overline{\mathcal{F}}_{G}(\nu)\) (cotransform of Fourier of ν). \ldots.. 471 \(\overline{\mathcal{F}}_{G}(f)\) (cotransform of Fourier of f). \ldots.. 471 \(\|u\|_{2}\) (norm in End(Eπ)). \ldots.. 471 \(\mathrm{F}^{2}(\widehat{G})\) (Hilbert sum of the End(Eπ)). \ldots.. 472 \(\mathcal{F}_{G}(f)\) (Fourier transform of f). \ldots.. 473 \(\mathcal{C}(\widehat{G})\) (functions on \(\widehat{G}\)). \ldots.. 475 \(M^{1}(G^{\sharp})\) (central bounded measures). \ldots.. 475 \(\mathrm{FS}(\pi)\) (Frobenius--Schur indicator). \ldots.. 476 \(M_{2k}(\varrho)\) (absolute moment of order 2k). \ldots.. 476 \(S^{2}(E)\) (symmetric square of an A-module). \ldots.. 478 \(\wedge^{2}E\) (exterior square of an A-module). \ldots.. 478

\specialchapter{TERMINOLOGICAL INDEX}

A Calkin algebra, 75 natural, 130 Fredholm alternative, 79 of Larsen, 476 Weakly mixing map,
322
Compact linear map, 2 Hilbert--Schmidt, 23 of finite rank, 4 of finite trace, 164 defined by a kernel, 25 nuclear, 169 strict, 2 Approximation space of ---, 14 property of ---, 14 Self-adjoint partial isometry ---, 238 essentially partial isometry ---, 238 Automorphic (representation ---), 405

B Basal (natural algebra ---), 131 Bottom (part --- of the spectrum), 298 Bicommutant (theorem of the ---), 326 Biregular (representation ---), 406 Bochner (theorem of ---), 455

Bounded (representation ---), 374 Bornification, 262 Unit ball, 1 Britton (lemma of ---), 506 C Functional calculus, 185 universally measurable, 272 Calkin (algebra of ---), 75 Character, 374 central, 390 Deficiency couple of ---, 256 index of ---, 256 Exterior square, 478 symmetric, 478 Cauchy (measure of ---), 515 Cauchy--Kowalewski (theorem of ---), 343 Central character ---, 390 measure ---, 402 Matrix coefficient, 386 of finite dimension, 386 diagonal, 386 Core of a partial operator, 231 Compact (linear map ---), 2 Completely continuous, 104, 107 Isotypic component, 394 Boundary condition, 258 Conilespace, 45

Conjugate (representation ---), 382 Conorm, 61 Cheeger constant, 323 Gelfand--Naimark--Segal construction, 433 Contragredient (representation ---), 377 Convolution of a distribution and a function, 339 Correspondence, 224 Fourier cotransform, 196, 471 of distributions, 219 Fourier cotransform, 471 Gelfand pair, 512 of Kazhdan, 499 Deficiency couple, 256 Weyl equidistribution criterion, 470 Cyclic Hilbertian realization ---, 433, 438, 443 vector ---, 191

D Polar decomposition, 291 Formal degree, 423 Partial derivative of a distribution, 208 Diagonal endomorphism ---, 24 endomorphism --- in a basis, 146 partial operator --- in a basis, 234 Direct morphism ---, 56 vector subspace ---, 55 Distribution, 203 with compact support, 336 tempered, 215 Domain of a partial operator, 224 Unitary dual, 393

E Fredholm endomorphism, 41 of multiplication, 186 of Riesz, 46

of finite trace, 23, 164 diagonal, 24 diagonal in the basis B, 146 Kazhdan set, 499 resolvent, 244 Schrödinger equation, 431 integral, 80 Equivariant (functions ---), 408 Approximation space, 14 Schwartz space, 209 Sobolev space, 222 Essential (spectrum ---), 297 Self-adjoint extension, 255 of Friedrichs, 295 almost analytic, 281 Extension of a partial operator, 226 F Sheaf, 204 Closure of a partial operator, 228 Vector bundle, 36 Faithful (representation ---), 374 Function with polynomial growth, 214 central, 402 of Loewner, 326 of Schwartz, 209 of negative type, 495 of positive type, 442 of interpolation, 327 representative, 386 spherical, 513 test, 200 Partial Hermitian form, 291 positive linear form, 436 partial positive, 291 associated partial positive, 291 closed partial positive, 292 invariant sesquilinear, 380 Formal (degree ---), 423 Formally symmetric (differential operator ---), 243 Helffer--Sjöstrand formula, 284 Leibniz formula, 196, 209

of Plancherel, 474 of Stone, 353 Trace formula, 488 Fredholm alternative of ---, 79 map of ---, 41 endomorphism of ---, 41 index of a map of ---, 43 Friedrichs (extension of ---), 295 Frobenius--Schur (indicator of ---), 476 G Gelfand--Naimark (theorem of ---), 442 Gelfand--Naimark--Segal (construction of ---), 433 Gelfand--Raikov (theorem of ---), 454 Infinitesimal generator, 428 Graph of a partial operator, 224 H Hamburger (moment problem of ---), 359 High (part --- of the spectrum), 298 Helffer--Sjöstrand (formula of ---), 284 Hellinger--Toeplitz (theorem of ---), 240 Hilbertian realization ---, 433, 438, 443 cyclic Hilbertian realization ---, 433, 438, 443 Hille--Yosida (theorem of ---), 346 I Identity (involutive automorphism ---), 478 μ-essential image, 181 Frobenius--Schur indicator, 476 Index, 43 of deficiency, 256 Kato inequality, 344 Irreducible matrix ---, 132 representation ---, 378 Isomorphism of representations, 375 Isotypic (component ---), 394 K Kato

inequality of ---, 344 Kato--Rellich (theorem of ---), 306 Kazhdan pair of ---, 499 set of ---, 499 property {]}T{[} of ---, 498 Kesten (theorem of ---), 323 Krein--Rutman (theorem of ---), 94

L Laplacian, 243 of Dirichlet, 310 Larsen (alternative of ---), 476 Leibniz (formula of ---), 196 Schur lemma, 12, 386 Lewy (theorem of ---), 343 Loewner (theorem of ---), 330 Weyl law, 311 Lomonosov (theorem of ---), 13

M Irreducible matrix, 132 primitive, 138 reducible, 132 Maxi-min, 154, 300 Mercer (theorem of ---), 177 Central measure, 402 of Cauchy, 515 spectral, 191, 268 tempered, 216 Mini-max, 154, 300 Bounded below (partial operator ---), 238 Absolute moment, 476 Moments of a measure, 359 moment problem of --- of Hamburger, 359 moment problem of --- of Stieltjes, 364 Moore--Penrose (pseudo-inverse of ---), 320 Morphism of representations, 375 direct, 56 Multiplication endomorphism of ---, 186 partial operator of ---, 231 Multiplicity, 245 of π in ρ, 398

spectral, 82 finite spectral, 82

N Natural (algebra ---), 130 Negative (function of negative type ---), 495 Negligible modulo H (set ---),
408
Nilespace, 45, 48 Nondegenerate (representation ---),
373
Normal (partial operator ---), 265 Relative norm, 306 Kernel of a partial operator, 225 universally positive, 433 Nuclear (linear map ---),
169
\lettergroup{0}

Operator with index, 41 with kernel, 25 compact, 2 of Fredholm, 41 of Hilbert--Schmidt, 23 of Markov, 322 of multiplication, 186 of finite rank, 4 of Riesz, 46 of Sturm--Liouville, 350 of finite trace, 23 of Volterra, 118 differential, 235 scalar differential, 235 transposed differential, 242 integral with kernel, 29

Partial operator, 224 with dense domain, 227 with compact resolvent, 309 adjoint, 235 self-adjoint, 238 bijective, 225 core of a ---, 231 of multiplication, 231 domain of a ---, 224 essentially self-adjoint, 238 extension of a ---, 226 closable, 227 closed, 227

closure of a ---, 228 graph of a ---, 224 image under a ---, 225 inverse image under a ---, 225 injective, 225 normal, 265 kernel of a ---, 225 positive, 238 product, 226 inverse, 226 reduction of a ---, 226 relatively bounded, 306 representing a positive form,
292
resolvent of a ---, 244 sum, 226 spectrum of a ---, 244 surjective, 225 symmetric, 238 Orthogonality (relations of ---), 424,
458
\lettergroup{P}

Sets equal up to a locally negligible set, 425 Perron--Frobenius (theorem of ---),

Peter--Weyl (theorem of ---), 462 Plancherel (formula of ---), 474 Sensitive point of the spectrum, 83 Orthogonal polynomials of the first kind, 361 of the second kind, 363 Positive function of --- type, 442 partial operator ---, 238 Positive (linear form ---), 436 Primitive (matrix ---), 138 Maxi-min principle, 154, 300 mini-max principle, 154, 300 Tensor product of unitary representations, 385 external tensor product of unitary representations, 385 Spectral projector, 54, 277 Approximation property, 14 {]}T{[} of Kazhdan, 498

true almost everywhere modulo H, 408 Pseudo-spectrum, 250 Q Quasi-inverse, 40 canonical, 48 Quasi-isomorphism, 41 R Hilbertian realization, 433, 438, 443 cyclic Hilbertian realization, 433, 438, 443 Reducible (matrix ---), 132 Reduction, 226 Regular representation --- right, 405 representation --- left, 405 Orthogonality relations, 424, 458 Γ-automorphic representation, 405 biregular, 406 bounded, 374 conjugate, 382 contragredient, 377 square-integrable, 420 square-integrable modulo the center, 420 of orthogonal type, 476 of quaternionic type, 476 of real type, 476 of symplectic type, 476 faithful, 374 induced, 416 irreducible, 378 isomorphism of ---s, 375 linear, 374 continuous linear, 374 morphism of ---s, 375 nondegenerate, 373 tensor product of ---s unitary, 385 tensor power, 385 quasi-regular, 486 quotient, 376 regular right, 405 regular left, 405 restriction of a --- to a subgroup, 375 sum of copies of a ---, 376

Hilbertian sum of ---s unitary, 384 sub---, 375 trivial, 374 unitary, 379 Resolvent, 244 partial operator with compact ---, 309 Restriction of a representation, 375 S Schauder (theorem of ---), 9 Schrödinger (equation of ---), 431 Schur (lemma of ---), 12, 386 Schwartz space of ---, 209 function of ---, 209 Continuous semigroup, 345 Semisimple (unitary representation ---), 391 Sensitive (point --- of the spectrum), 83, 297 Smale (theorem of ---), 124 Sobolev (space of ---), 222 Summation by parts, 160 Subspace adapted to u, 153, 299 of coordinates, 132 invariant, 12 direct vector, 55 Subrepresentation, 375 generated by a subset, 376 Spectral measure ---, 268 multiplicity ---, 82 finite multiplicity ---, 82 projector ---, 277 theorem ---, 194, 266 Spectrum, 244 complex, 54 continuous, 245 essential, 83, 297 sensitive point of the ---, 83 residual, 245 sensitive, 297 stable, 86 Stieltjes (moment problem of ---), 364 Stone formula of ---, 353

theorem of ---, 428 Sequence of eigenvalues, 152 Support of a distribution, 336 Symbol, 127 Symmetric (partial operator ---), 238

\lettergroup{T}

Theorem of Bochner, 455 of Borsuk--Ulam, 64 of Cauchy--Kowalewski, 343 of Gelfand--Naimark, 442 of Gelfand--Raikov, 454 of Hellinger--Toeplitz, 240 of Hille--Yosida, 346 of Kato--Rellich, 306 of Kesten, 323 of Krein, Krasnoselskii and Milman, 64 of Krein--Rutman, 94 of Lewy, 343 of Loewner, 330 of Lomonosov, 13 of Mercer, 177 of Perron--Frobenius, 101 of Peter--Weyl, 462 of elliptic regularity, 341 of Rellich--Kondrachov, 335 of Schauder, 9 of Smale, 124 of Stone, 428 of von Neumann, 258 of the bicommutant, 326 of the Schwartz kernel, 342 spectral, 194, 266 Toeplitz (operator of ---), 127 Weak topology, 6 of the C \(^{r}\) -compact convergence, 34 weak, 449 T-resolvent, 285 Trace

finite trace map, 164 trace formula, 488

Fourier transformation, 196, 473 Fourier transformation of distributions, 219 Gelfand transformation, 190

Fourier transform, 473 T-resolvent (topology ---), 285 Trivial (representation ---), 374

U

Natural unital algebra, 131

Unitary dual ---, 393 tensor product of --- representations, 385 representation ---, 379 semisimple representation ---, 391 Hilbertian sum of representations ---, 384

Approximate unit, 437

Universally measurable functional calculus ---, 272 function ---, 183

V

Cauchy principal value, 337 Eigenvalue, 245 Eigenvalues (family of ---), 146,
234
Singular values, 151, 156 Cyclic vector, 191, 376 G-finite, 376 Volterra (operator of ---), 118 von Neumann theorem of ---, 258 bicommutant theorem of ---,
326
\lettergroup{W}

Weyl equidistribution criterion of ---, 470 inequalities of ---, 162 law of ---, 311

\specialchapter{SELF-ADJOINTNESS CRITERIA}

We summarize here the criteria stated in the text that allow one to prove that certain partial operators are self-adjoint.

We fix a complex Hilbert space E and a partial operator u on E. This operator is self-adjoint in each of the following cases:

\begin{enumerate}
\def\labelenumi{(\roman{enumi})}
\item
  There exists \(v \in \mathcal{L}(\mathrm{E})\) injective and Hermitian such that \(u = v^{-1}\) . (Proposition 9 of IV, p. 239);
\item
  The operator \(u\) is symmetric and \(\mathrm{dom}(u) = \mathrm{E}\) (Proposition 10 of IV, p. 240);
\item
  The operator \(u\) is symmetric and induces a bijection of \(\mathrm{dom}(u)\) onto E (corollary of Proposition 10 of IV, p. 240);
\item
  There exists \(\lambda \in \mathbf{C}\) such that \(u + \lambda 1_{\mathrm{E}}\) and \(u + \overline{\lambda} 1_{\mathrm{E}}\) are surjective (Proposition 11 of IV, p. 240);
\item
  There exists a partial operator \(v\) , closed with dense domain, such that \(u = v^{*} \circ v\) (Proposition 12 of IV, p. 241);
\item
  The operator \(u\) is symmetric, closed, and its deficiency index is \((0,0)\) (corollary of Proposition 26 of IV, p. 257), that is, \(\mathrm{Ker}(u^{*} - i) = \mathrm{Ker}(u^{*} + i) = \{0\}\) ;
\item
  The space E is of the form \(\mathrm{L}^{2}(\mathrm{X},\mu)\) , where X is a locally compact topological space and \(\mu\) is a measure on X, and u is the operator of multiplication by a function \(\mu\) -measurable on X which is locally \(\mu\) -almost everywhere real-valued (corollary of Proposition 23);
\item
  There exists a normal partial operator v on E and a universally measurable function \(f \in \mathcal{L}_{\mathrm{u}}(\mathrm{Sp}(v))\) with real values such that \(u = f(v)\) (Corollary 1 of IV, p. 273);
\item
  There exists a partial Hermitian form q on E such that u is the partial operator representing q (Proposition 14 of IV, p. 293);
\item
  One has \(u = v + w\) where v is a self-adjoint partial operator and w is a symmetric partial operator bounded relative to v satisfying \(\|w\|_{v} < 1\) (Theorem 5 of IV, p. 306);
\item
  The operator \(u\) is the infinitesimal generator of a unitary representation \(\varrho\) of \(\mathbf{R}\) in E (Theorem 1 of V, p. 428).
\end{enumerate}
